% Beispiel für das Corporate Design der TU Chemnitz (Stand 2019)
\DocumentMetadata{}
\documentclass{ltx-talk}

\usepackage{ltxtalk-tuc-2019}
\settucthemelogo{../../images/logos/tuc_white.pdf}
\settucthemetitlelogo{../../images/logos/osg.png}
\settucthemeurl{https://www.tu-chemnitz.de/}
\useltxtalktheme{tuc-2019}

\title{Echtzeitsysteme}
\subtitle{0. Kapitel\\Organisatorisches}
\author{Prof. Matthias Werner}
\institute{Professur Betriebssysteme}
\date{Sommersemester 2026}

\begin{document}

\maketuctitle

\section{Organisatorisches}
\subsection{Willkommen}

\begin{frame}
  \frametitle{Willkommen zur Lehrveranstaltung}
  \framesubtitle{Echtzeitsysteme}

  Das Theme bildet die wesentlichen Elemente des TU-Chemnitz-Designs ab:
  \begin{itemize}
    \item grünes Kopfband mit Universitätslogo und Veranstaltungskontext,
    \item grüne Folientitel und Aufzählungszeichen,
    \item dreigeteilte Fußzeile mit Autor, Foliennummer und Webadresse.
  \end{itemize}
\end{frame}

\subsection{Regularien}

\begin{frame}
  \frametitle{0.1 Regularien}
  \begin{itemize}
    \item Alle Informationen und Materialien werden über die Lernplattform
      bereitgestellt.
    \item Fragen können in der Vorlesung und in den Übungen gestellt werden.
    \item Die Beispiele dieser Präsentation sind vollständig mit
      \texttt{ltx-talk} gesetzt.
  \end{itemize}

  \begin{block}{Hinweis}
    Veranstaltungsbezogene Bilder und Inhalte gehören zur Präsentation, nicht
    zum Theme.
  \end{block}
\end{frame}

\section{Inhalt}
\subsection{Überblick}

\begin{frame}
  \frametitle{Unsere Reise}
  \begin{enumerate}
    \item Begriffe und Konzepte
    \item Anforderungen und Spezifikationen
    \item Prioritäts- und zeitgesteuertes Scheduling
    \item Weiche und harte Echtzeit
  \end{enumerate}
\end{frame}

\end{document}
